% Part: lambda-calculus
% Chapter: introduction
% Section: syntax

\documentclass[../../../include/open-logic-section]{subfiles}

\begin{document}

\olfileid{lam}{int}{syn}
\olsection{λ-演算的语法}

我们从变元序列 $x$, $y$, $z$,~\dots 以及一些常项符号 $a$, $b$, $c$,~\dots 出发。项的集合归纳定义如下：
\begin{enumerate}
\item 每个变元都是项。
\item 每个常项都是项。
\item 如果 $M$ 和 $N$ 是项，那么 $(MN)$ 也是项。
\item 如果 $M$ 是项且 $x$ 是变元，那么 $(\lambd[x][M])$ 也是项。
\end{enumerate}

形如 $(MN)$ 的项称为\emph{应用}，形如 $(\lambd[x][M])$ 的项称为\emph{抽象}。

不含任何常项的系统称为\emph{纯} $\lambd$-演算。我们主要在纯 $\lambd$-演算中进行讨论，因此所有小写字母都表示变元。我们用大写字母（$M$, $N$ 等）来表示 $\lambd$-演算中的项。

我们将遵循以下记法约定：

\begin{conv}
\begin{enumerate}
\item 省略括号时，应用的结合从左向右进行。例如，如果 $M$, $N$, $P$, $Q$ 是项，那么 $MNPQ$ 是 $(((MN)P)Q)$ 的缩写。
\item 同样，省略括号时，$\lambd$-抽象应取得尽可能广的辖域。例如，$\lambd[x][MNP]$ 读作 $(\lambd[x][((MN)P)])$。
\item 一个 $\lambd$ 可用于抽象多个变元。例如，$\lambd[xyz][M]$ 是 $\lambd[x][\lambd[y][\lambd[z][M]]]$ 的缩写。
\end{enumerate}
\end{conv}

例如，
\[
\lambd[xy][xxyx \lambd[z][xz]]
\]
是
\[
\lambd[x][\lambd[y][((((xx)y)x)(\lambd[z][(xz)]))]].
\]
的缩写。
你应当记住这些约定。一开始它们可能会令你抓狂，但你会习惯的；过一段时间，比起处理一大堆括号，这些约定带来的困扰就要小得多了。

仅约束变元的名称不同的两个项称为 $\alpha$-等价的；例如 $\lambd[x][x]$ 和 $\lambd[y][y]$。为方便起见，可以将这些项视为“同一个”项；换言之，当我们说 $M$ 和 $N$ 相同时，也意味着“至多相差约束变元的重命名”。处于 $\lambd$ 辖域内的变元称为“约束的”，其余的变元称为“自由的”。前面的例子中没有自由变元；但在
\[
(\lambd[z][yz])x
\]
中，$y$ 和 $x$ 是自由的，而 $z$ 是约束的。

\end{document}